% !TEX program = lualatex
% !TEX options = -synctex=1 -interaction=nonstopmode -file-line-error -shell-escape -halt-on-error %DOC%
% !TEX encoding = UTF-8 Unicode

\documentclass[twoside, a4paper]{extreport}
\usepackage{scrextend}
\changefontsizes{16pt} % ustawia 16pt jako bazę dla całego dokumentu
% -----------------------------------------------------------------------------------
% ------------------------- USTAWIENIA STRONY--------------------------
% -----------------------------------------------------------------------------------
\usepackage{geometry}
\geometry{
	left=10mm,
	right=10mm,
	top=15mm,
	bottom=15mm,
	marginparwidth=0mm,
	marginparsep=0mm
	}
\newcommand{\wstawplik}[1]{\input{../include/#1}}%skraca kod
\wstawplik{headerofficium-kopia}
% elementy zmienne co do headerofficium:
\graphicspath{{../img/}} % katalog obrazów
%\svgpath{{../include/}}
\gresetlastline{ragged}
% koniec 

% debugowanie:
% %\iffalse %*** %\fi %***
% \tracingall przed problematycznym kodem - bardzo dużo informacji
% lualatex+se+debug - dodatkowe opcje: -file-line-error i -halt-on-error 
% \errorcontextlines
\def\mydebug{1}% 0 - debugowanie, 1 - bez debugowania (wyłącza bloki \if\mydebug1)
\begin{document}
% -----------------------------------------------------------------------------------
% ------------------------- STRONA TYTUŁOWA --------------------------
\thispagestyle{empty}
\BeginParallel{0.5}
\fontsize{17}{16}\selectfont
\latN{Kýrie, eléison. Chríste, eléison. Kýrie, eléison.}\polN{Kyrie, eleison. Chryste, eleison. Kýrie, eléison.}%
\latN{Chríste, áudi nos. Chríste, exáudi nos.}\polN{Chryste, usłysz nas. Chryste, wysłuchaj nas.}%
\latN{Páter de caélis, Déus, \textbf{miserére nóbis.}}\polN{Ojcze z nieba, Boże, \textbf{zmiłuj się nad nami.}}%
\latN{Fíli, Redémptor múndi, Déus, }\polN{Synu, Odkupicielu świata, Boże,}%
\latN{Spíritus Sáncte, Déus, }\polN{Duchu Święty, Boże,}%
\latN{Sáncta Trínitas, únus Déus, }\polN{Święta Trójco, Jedyny Boże,}%
\latN{Sáncta María, \textbf{ora pro nobis.}}\polN{Święta Maryjo, \textbf{módl się za nami.}}%
\latN{Sáncte Ióseph,}\polN{Święty Józefie,}%
\latN{Próles Dávid ínclita,}\polN{Przesławny potomku Dawida,}%
\latN{Lúmen Patriarchárum,}\polN{Światło Patriarchów,}%
\latN{Déi Genitrícis Spónse,}\polN{Oblubieńcze Bogarodzicy,}%
\latN{Cústos pudíce Vírginis,}\polN{Przeczysty stróżu Dziewicy,}%
\latN{Fílii Déi nutrície,}\polN{Żywicielu Syna Bożego,}%
\latN{Chrísti defénsor sédule,}\pol{Troskliwy obrońco Chrystusa,}%
\latN{Almae Famíliae praéses,}\polN{Głowo Najświętszej Rodziny,}%
\latN{Ióseph iustíssime,}\polN{Józefie najsprawiedliwszy, }%
\latN{Ióseph castíssime,}\polN{Józefie najczystszy,}%
\latN{Ióseph prudentíssime,}\polN{Józefie najroztropniejszy,}%
\latN{Ióseph fortíssime,}\polN{Józefie najmężniejszy,}%
\latN{Ióseph oboedientíssime,}\polN{Józefie najposłuszniejszy,}%
\latN{Ióseph fidelíssime,}\polN{Józefie najwierniejszy,}%
\latN{Spéculum patiéntiae,}\polN{Zwierciadło cierpliwości,}%
\latN{Amátor paupertátis,}\polN{Miłośniku ubóstwa,}%
\latN{Exémplar opíficum,}\polN{Wzorze pracujących,}%
\latN{Domésticae vítae décus,}\polN{Ozdobo życia rodzinnego,}%
\latN{Cústos vírginum,}\polN{Opiekunie dziewic,}%
\latN{Familiárum colúmen,}\polN{Podporo rodzin,}%
\latN{Solátium miserórum,}\polN{Pociecho nieszczęśliwych,}%
\latN{Spés aegrotántium,}\polN{Nadziejo chorych,}%
\latN{Patróne moriéntium,}\polN{Patronie umierających,}%
\latN{Térror daémonum,}\polN{Postrachu duchów piekielnych,}%
\latN{Protéctor sánctae Ecclésiae,}\polN{Opiekunie Kościoła świętego,}%
\latN{Agnus Déi, quí tóllis peccáta múndi, párce nóbis, Dómine.}\polN{Baranku Boży, który gładzisz grzechy świata, przepuść nam, Panie.}%
\latN{Agnus Déi, quí tóllis peccáta múndi, exáudi nos, Dómine.}\polN{Baranku Boży, który gładzisz grzechy świata, wysłuchaj nas, Panie.}%
\latN{Agnus Déi, quí tóllis peccáta múndi, miserére nóbis.}\polN{Baranku Boży, który gładzisz grzechy świata, zmiłuj się nad nami.}%
\latN{\red{\Vbar.} Constítuit éum dóminum dómus súae.}\polN{\red{\Vbar.} Ustanowił go panem domu swego.}%
\latN{\red{\Rbar.} Et príncipem ómnis possessiónis súae.}\polN{\red{\Rbar.} I zarządcą wszystkich posiadłości swoich.}%
\EndParallel
\end{document}